\section{Reproducibility and software citation}
\label{sec:reproducibility}

\subsection{Reproducing the results in this paper}

Every number, table and figure in Sections~\ref{sec:validation}
and~\ref{sec:examples} is produced by scripts in
\code{paper/phynetpy-1.0.0/}. Running

\begin{lstlisting}[language=bash, basicstyle=\ttfamily\small]
export PHYLONET_JAR=/path/to/phylonet.jar
python paper/phynetpy-1.0.0/benchmarks/run_all.py
python paper/phynetpy-1.0.0/benchmarks/make_analyses_table.py
python paper/phynetpy-1.0.0/listings/paper_listings.py
\end{lstlisting}

regenerates the benchmark JSON, rewrites \code{tables/*.tex} and
\code{figures/performance.pdf} from it, rebuilds Table~\ref{tab:analyses} from
the live dispatch registry, and executes every code listing. The PhyloNet
comparison is skipped with a warning if no jar is configured; the rest needs
nothing beyond the package and its dependencies. The tables under
\code{tables/} are generated files and carry a header saying so.

\subsection{Environment}

Measurements in this paper were taken under the following configuration, which
each benchmark script also records in the \code{meta} block of its own JSON
output:

\begin{itemize}
  \item \phynetpy\ 0.6.0.\footnote{[TODO: re-run every benchmark against the
    tagged 1.0.0 release and update this section, the abstract and
    Section~\ref{sec:validation}. The measurements reported here were taken on
    the 0.6.0 codebase, which is 1.0.0 without the two unbuilt cells of
    Table~\ref{tab:analyses}.]}
  \item CPython 3.14.2 on Windows 11 (build 10.0.26200), AMD64.
  \item AMD Family 25 Model 97 (Zen 4 class), single-threaded unless stated.
  \item NumPy 2.3.5, SciPy 1.16.3, NetworkX 3.6.1.
  \item PhyloNet 3.8.2 (\code{phylonet.jar}, 41{,}019{,}851 bytes) under
    Java 1.8.0\_491, with BEAGLE for the Felsenstein cross-check.
\end{itemize}

Random seeds are fixed in every script and recorded per measurement. The
head-to-head comparison runs three replicates per condition, replicate $i$ using
seed $20260925 + i$ for both simulation and \phynetpy's search; the component
benchmarks use seeds 7 and 11, and the marker timings use seed 1. PhyloNet's
search is stochastic and its command-line interface is not seeded here, so its
variation across replicates is uncontrolled -- which is part of why the
comparison is replicated rather than run once. The generating network for each
replicate is stored as Newick in the benchmark JSON alongside its results,
together with the gene-tree concordance that replicate produced, so any trial can
be reconstructed exactly.

\paragraph{Timing protocol.} Microbenchmarks report the median of repeated
batches measured with \code{time.perf_counter}, after one warm-up call; loop and
repeat counts are stored per measurement. Interpreter start-up and module import
are excluded and recorded separately. For the PhyloNet comparison, PhyloNet's
self-reported running time is used because it excludes JVM start-up and is
therefore comparable to \phynetpy's in-process timing; total wall time including
JVM start-up is also recorded in the JSON. No result in this paper uses
multi-threading or GPU offload.

\paragraph{Datasets.} All data are simulated by \phynetpy\ itself, from seeds
recorded in the scripts, so no external archive is required. No empirical dataset
is analysed in this paper.

\subsection{Citing \phynetpy}

\begin{itemize}
  \item Software: \phynetpy\ version 1.0.0.
  \item Repository: \url{https://github.com/NakhlehLab/PhyNetPy}, tag
    \code{v1.0.0}.\footnote{[TODO: create the \code{v1.0.0} tag and confirm this
    reference before submission; the current release tag is \code{v0.6.0}.]}
  \item Package: \url{https://pypi.org/project/phynetpy/}.
  \item License: MIT.
  \item Archived release: \cite{TODO_PhyNetPy_archive}.
\end{itemize}

Methods reimplemented from published work should be cited to their original
sources rather than to \phynetpy\ or to PhyloNet generically; the relevant
citation for each is given in Table~\ref{tab:analyses}.
